\begin{description}
\item[Q1.\ Tightness of Theorem 2 vs.\ PGM.]  We observe empirical PGM
  $\log_2(\mathrm{err})$-slopes of $\{0.71, 0.56, 0.45\}$ on
  $S_3, D_4, \mathbb{Z}_2^{3}$, all comparable to the paper's guaranteed slope
  of $\tfrac12$, but with the striking pattern that the slope \emph{decreases}
  as $r$ (the number of subgroups) grows.  Is this a genuine
  information-theoretic feature (many-way discrimination becomes intrinsically
  harder) or an artefact of PGM sub-optimality at small $s$?  Fitting the
  paper's own \textsf{Test} on the same instances would answer this directly.

\item[Q2.\ Extrapolated $s^*/\log_2|G|$ constant across group flavour.]  Our
  three groups have wildly different structure but $s^*/\log_2|G|$ lands in
  the tight window $[3.9, 5.3]$.  Is this constant sensitive to
  \emph{class-2} vs.\ perfect group structure?  A follow-up sweep across
  $S_n$ for $n=3,4,5$ (where subgroup counts grow super-polynomially in
  $\log|G|$) would test whether the constant is really $\Theta(1)$ or
  quietly hides a $\log r / \log|G|$ factor.

\item[Q3.\ Coset-state distinguishability on non-Abelian $G$ where HSP is
  known hard.]  The dihedral group HSP is a benchmark hardness case
  (Kuperberg's subexponential algorithm).  Our $D_4$ result shows
  information-theoretically the coset states already distinguish subgroups
  after $s\sim 8$ queries, yet no efficient quantum algorithm is known.
  Where in the algorithmic pipeline does the exponential-time barrier
  actually sit --- Fourier decomposition of $\rho_H^{\otimes s}$ into
  irreps, or the classical postprocessing of PGM POVM elements?  A concrete
  next step: measure the \emph{circuit depth} required to synthesise our
  numerical PGM POVM to precision $\epsilon$ on $D_4$ at $s=4$.

\item[Q4.\ The paper's Neumann-series bound
  $\|\Delta\|_\infty \le r^{-2}$.]  Theorem's exact algorithm needs
  $s = \lceil 2\log(4r^3)\rceil \approx 12$--$14$ for our groups, well past
  our tractable regime.  Our tables show that $\|\Delta\|_\infty$ actually
  becomes small far earlier: at $s=4$ the max off-diagonal
  $\Pr[\mu\ne H|H]$ is already $\le 0.03$ for $S_3$, i.e.\ $\|\Delta\|_\infty
  \le 0.15 \ll r^{-2} = 0.028$ (just barely).  Is $s = O(\log r)$ (rather than
  $s = O(\log r^3)$) actually enough for the exact algorithm, and does the
  paper leave a factor-3 in the exponent on the table?

\item[Q5.\ Priors mismatch on \emph{semantic} subgroup families.]  We used a
  uniform prior over subgroups, matching the paper's setup.  In real
  applications --- period-finding, graph isomorphism --- the effective prior
  is highly non-uniform (e.g.\ cyclic subgroups dominate for Shor).
  How much does $\min_H \Pr[H|H]$ improve if we restrict the ensemble to a
  \emph{physically motivated} family (say, all cyclic subgroups of $D_n$)
  and re-run PGM?  This would test whether the paper's $4r$ factor can be
  refined to $4|\mathcal{K}|$ for the relevant subgroup family
  $\mathcal{K}$, which would substantially tighten Thm.~2 for structured
  HSP.
\end{description}
